\subsection*{Theories}
\label{sec:axiom-systems-theories}

\begin{tcolorbox}[colback=structuretheorybox, colframe=structuretheoryborder, arc=2pt,
  left=8pt, right=8pt, top=6pt, bottom=6pt,
  title={\small\textbf{Theory Presentation: Arithmetic Consequences}},
  fonttitle=\small\bfseries]
Start with the arithmetic-style axiom system from the previous section. It
uses the arithmetic signature and includes assertions such as
\[
  \forall x\;(x+0=x),
\]
\[
  \forall x\;(S(x)\neq 0),
\]
\[
  \forall x\forall y\;
  (S(x)=S(y)\rightarrow x=y).
\]

Those axioms are the starting assertions. From them, together with other
arithmetic axioms and ordinary logical reasoning, additional statements may be
established. For instance, one may later prove new equations, cancellation
facts, uniqueness facts, and order facts that were not listed as axioms.

A theory contains the original axioms. A theory also contains the statements
derivable from those axioms. Thus the axioms are the generators, while the
theory is the full collection of what those generators imply.
\end{tcolorbox}

\begin{remark*}[Structural remarks]
\textbf{Axioms vs theory.}
Axioms are starting assertions. A theory includes all logical consequences of
those assertions. The distinction is the next step after the vocabulary and
assertion distinction: a signature supplies vocabulary, axioms supply adopted
assertions, and a theory contains everything derivable from those assertions.

\textbf{Finite axioms, infinite theory.}
A small axiom system can generate an enormous collection of consequences. The
listed axioms may be finite, but the statements derivable from them need not
be. This is why a theory is usually larger than the axiom system that presents
it.

\textbf{Forward roadmap.}
The next conceptual step is semantic:
\[
\begin{array}{c}
\text{Theory}\\
\Downarrow\\
\text{Models}
\end{array}
\]
A model will later be introduced as a mathematical structure in which the
statements of a theory hold.
\end{remark*}

\begin{definitionbox}{Definition (Theory)}
\begin{definition}[Theory]
\label{def:axiom-systems-theory}
A \emph{theory} over a first-order language $\mathcal{L}$ is a collection of
closed sentences in that language:
\[
  T\subseteq\operatorname{Sentences}(\mathcal{L}).
\]
If $\Gamma$ is a chosen set of axioms, the theory generated by $\Gamma$ is the
deductive closure
\[
  T=\{\varphi\in\operatorname{Sentences}(\mathcal{L})\mid
      \Gamma\vdash\varphi\}.
\]
\end{definition}
\end{definitionbox}

\begin{remark*}[Standard quantified statement]
Syntactically, a sentence $\varphi$ belongs to the theory generated by
$\Gamma$ exactly when $\varphi$ is derivable from $\Gamma$. Semantically, for
an $\mathcal{L}$-structure $\mathfrak{M}$, the complete theory of
$\mathfrak{M}$ is
\[
  \operatorname{Th}(\mathfrak{M})
  =\{\varphi\in\operatorname{Sentences}(\mathcal{L})\mid
      \mathfrak{M}\vDash\varphi\}.
\]
For every sentence $\varphi$, either $\mathfrak{M}\vDash\varphi$ or
$\mathfrak{M}\vDash\neg\varphi$, so $\operatorname{Th}(\mathfrak{M})$
decides every sentence of the language.
\end{remark*}

\begin{remark*}[Interpretation]
A theory is not merely the list of axioms. It is the closure of those axioms
under derivability: every statement that follows from the axioms belongs to
the theory. The axioms start the system; the theory records everything the
system proves.
\end{remark*}

\begin{dependencies}
\begin{itemize}
  \item \hyperref[def:language]{Language}
  \item \hyperref[def:axiom-system]{Axiom System}
\end{itemize}
\end{dependencies}

\begin{definitionbox}{Definition (Theorem)}
\begin{definition}[Theorem]
\label{def:axiom-systems-theorem}
A \emph{theorem} of a theory $T$ is a sentence $\varphi$ for which there is a
finite formal derivation from $T$, written
\[
  T\vdash\varphi.
\]
Equivalently, in proof-theoretic terms, $\varphi$ is a theorem of $T$ exactly
when some derivation has leaves drawn from logical axioms or sentences of
$T$, uses valid inference rules, and has $\varphi$ as its root.
\end{definition}
\end{definitionbox}
\begin{remark*}[Standard quantified statement]
Syntactically,
\[
  \operatorname{Theorem}(T,\varphi)
  \iff \exists d\ (d:T\vdash\varphi).
\]
Semantically, $\varphi$ is entailed by $T$ when it is true in every model of
$T$:
\[
  T\vDash\varphi
  \iff
  \forall\mathfrak{M}\,
  \bigl(\mathfrak{M}\vDash T \implies \mathfrak{M}\vDash\varphi\bigr).
\]
For first-order logic, the completeness theorem connects these readings:
$T\vdash\varphi$ exactly when $T\vDash\varphi$.
\end{remark*}

\begin{remark*}[Predicate reading]
The predicate form records the relation or property named by the statement in
the displayed logical context.
\end{remark*}


\begin{remark*}[Interpretation]
Within an axiom system, a theorem is not an additional starting assumption. It
is a statement reached from the adopted axioms by reasoning. Axioms are given;
theorems are obtained.
\end{remark*}

\begin{dependencies}
\begin{itemize}
  \item \hyperref[def:axiom-systems-theory]{Theory}
\end{itemize}
\end{dependencies}

\begin{definitionbox}{Definition (Consequence)}
\begin{definition}[Consequence]
\label{def:axiom-systems-consequence}
A sentence $\varphi$ is a \emph{semantic consequence} of a set of premises
$\Gamma$, written $\Gamma\vDash\varphi$, if every model of $\Gamma$ is also a
model of $\varphi$. Equivalently, there is no structure in which all sentences
of $\Gamma$ are true while $\varphi$ is false.
\end{definition}
\end{definitionbox}
\begin{remark*}[Standard quantified statement]
Semantic consequence is the universal condition
\[
  \Gamma\vDash\varphi
  \iff
  \forall\mathfrak{M}\,
  \Bigl(
    (\forall\gamma\in\Gamma,\ \mathfrak{M}\vDash\gamma)
    \implies
    \mathfrak{M}\vDash\varphi
  \Bigr).
\]
Syntactic consequence, written $\Gamma\vdash\varphi$, means that there exists
a finite valid deduction of $\varphi$ from $\Gamma$. For first-order logic,
the completeness theorem says that $\Gamma\vDash\varphi$ exactly when
$\Gamma\vdash\varphi$.
\end{remark*}

\begin{remark*}[Predicate reading]
The predicate form records the relation or property named by the statement in
the displayed logical context.
\end{remark*}


\begin{remark*}[Interpretation]
Consequence is the relation that connects a theory to what it implies. If the
axioms generate a theory, then consequences are the statements carried along
by those axioms and by the statements already obtained from them.
\end{remark*}

\begin{dependencies}
\begin{itemize}
  \item \hyperref[def:axiom-systems-theory]{Theory}
\end{itemize}
\end{dependencies}

\begin{definitionbox}{Definition (Deductive Closure)}
\begin{definition}[Deductive Closure]
\label{def:deductive-closure}
Let $\Gamma\subseteq\operatorname{Sentences}(\mathcal{L})$ be a set of
premises. The \emph{deductive closure} of $\Gamma$, denoted
$\operatorname{Cn}(\Gamma)$, is the set of all sentences derivable from
$\Gamma$:
\[
  \operatorname{Cn}(\Gamma)
  =
  \{\varphi\in\operatorname{Sentences}(\mathcal{L})\mid
    \Gamma\vdash\varphi\}.
\]
\end{definition}
\end{definitionbox}

\begin{remark*}[Standard quantified statement]
A set of sentences $T$ is deductively closed exactly when
\[
  T=\operatorname{Cn}(T),
\]
or equivalently, for every sentence $\varphi$,
\[
  T\vdash\varphi \implies \varphi\in T.
\]
The closure operator satisfies the standard Tarskian closure laws:
\[
  \Gamma\subseteq\operatorname{Cn}(\Gamma),
  \qquad
  \Gamma\subseteq\Delta\implies
  \operatorname{Cn}(\Gamma)\subseteq\operatorname{Cn}(\Delta),
  \qquad
  \operatorname{Cn}(\operatorname{Cn}(\Gamma))
  =\operatorname{Cn}(\Gamma).
\]
By completeness for first-order logic, this syntactic closure agrees with the
semantic closure
\[
  \{\varphi\in\operatorname{Sentences}(\mathcal{L})\mid
    \Gamma\vDash\varphi\}.
\]
\end{remark*}

\begin{remark*}[Interpretation]
Deductive closure is the operation that turns an axiom system into the full
theory it generates. Starting with the axioms, close under derivability; the
result is the theory determined by those starting assertions.
\end{remark*}

\begin{dependencies}
\begin{itemize}
  \item \hyperref[def:axiom-system]{Axiom System}
\end{itemize}
\end{dependencies}
